\SourceLeaf{313}
Da nun die Nenner in II. und III. beide denselben Wert, nämlich

\[
\tfrac{1}{2}\rho\left\{\sin(2A-2A')+\sin(2A'-2A'')+\sin(2A''-2A)\right\},
\]

haben, ergibt sich aus diesen Gleichungen durch Division:

\[
\operatorname{tg}\psi = -\dfrac{mmh\cos A+m'm'h'\cos A'+m''m''h''\cos A''}{mmh\sin A+m'm'h'\sin A'+m''m''h''\sin A''},
\]

woraus zwei Werte für $\psi$ und damit auch für $R$ folgen.]

\bigskip
\begin{center}
\hrulefill
\end{center}
\bigskip

\Needspace{9\baselineskip}
\begin{center}
BEMERKUNG.
\end{center}

Die vorstehende Notiz, die sich in einem Handbuch findet und vermutlich aus dem Jahr 1816 stammt, ist sehr flüchtig geschrieben, sodass zahlreiche Schreibfehler in den Formeln berichtigt werden mussten. Sie bricht bei Formel II der zweiten Lösungsmethode ab. Dabei mussten der Nenner und der Faktor $\Delta$ ergänzt werden, was im Druck nicht kenntlich gemacht werden konnte. Zur Verdeutlichung wurden dem Text zwei Figuren hinzugefügt.

Es ist anzunehmen, dass \textsc{Gauss} seine Lösungen gefunden hat, indem er von der Darstellung der Ecken des Dreiecks durch komplexe Größen ausging, denn dabei ergeben sich genau die Formeln des Textes, wenn man die Gleichungspaare, die den Kosinus und den Sinus desselben Winkels enthalten, jeweils mithilfe des Symbols $i$ zu einer einzigen Gleichung zusammenfasst. Dass \textsc{Gauss} die Darstellung der Punkte einer Ebene durch komplexe Größen schon früh bei geometrischen Aufgaben verwendet hat, zeigt zum Beispiel die auf der folgenden Seite abgedruckte Notiz über die Konstruktion der Aufgabe von \textsc{Pothenot} sowie eine handschriftliche Bemerkung zu \textsc{Schumachers} Übersetzung der \emph{Geometrie der Stellung} von \textsc{Carnot} aus dem Jahr 1810, die sich ebenfalls auf die Geometrie des Dreiecks bezieht und bereits in Bd.\ IV.\ S.\ 396 bis 397 erschienen ist.

\hfill \textsc{Stäckel}.

\bigskip
\begin{center}
\hrulefill
\end{center}
